% Part: propositional-logic

\documentclass[../../include/open-logic-part]{subfiles}

\begin{document}

\olpart{pl}{Logica van samengestelde beweringen (propositielogica)}

\begin{editorial}
  Dit deel gaat over klassieke logica van samengestelde beweringen,
  oftewel propositielogica. Het eerste hoofdstuk behandelt vooral de
  basis en zet alleen definities en resultaten op een rij. Veel bewijzen
  staan er niet bij; de lezer werkt die als opgaven uit. De stukken over
  bewijssystemen (stelsels van bewijsregels) en de volledigheidsstelling
  komen uit het deel over logica van de eerste orde. Bij het overnemen is
  de tag ``FOL'' op onwaar gezet. Daardoor blijven alle stukken over
  predicaten, termen en kwantoren weg. Waar die tekst over
  structuren~$\Struct{M}$ gaat, gaat deze versie over waardetoekenningen
  aan propositieletters (valuaties)~$\pAssign{v}$.

  Het plan is om dit deel verder uit te werken. Er komen ook meer
  onderwerpen en resultaten bij, zoals functionele volledigheid. Dat
  betekent dat elke formule kan worden vervangen door een equivalente
  formule die alleen de gekozen logische verbindingen gebruikt.
\end{editorial}

\olimport[syntax-and-semantics]{syntax-and-semantics}

\tagfalse{FOL}

\olimport[../first-order-logic/proof-systems]{proof-systems}

\iftag{prfSC}{%
  \olimport[../first-order-logic/sequent-calculus]{sequent-calculus}
}{}

\iftag{prfND}{%
  \olimport[../first-order-logic/natural-deduction]{natural-deduction}
}{}

\iftag{prfTab}{%
  \olimport[../first-order-logic/tableaux]{tableaux}
}{}

\iftag{prfAX}{%
  \olimport[../first-order-logic/axiomatic-deduction]{axiomatic-deduction}
}{}

\olimport[../first-order-logic/completeness]{completeness}

\tagtrue{FOL}

\OLEndPartHook
\end{document}
